% Generated by Listas/make_statements.py from pratica-lecture-01.md, item 21. Do not edit.
\textbf{Source.} Mas-Colell, Whinston \& Green (1995), \emph{Microeconomic Theory}, Ch. 8 Exercises, Exercise 8.B.5\^{}C, p.~262 {[}PDF 142{]}. Reproduced verbatim. Answer in the source: none in the book (see the MWG solutions manual in \texttt{Livros/}).



\begin{quote}

Consider the Cournot duopoly model (discussed extensively in Chapter 12) in which two firms, 1 and 2, simultaneously choose the quantities they will sell on the market, \(q_1\) and \(q_2\). The price each receives for each unit given these quantities is \(P(q_1, q_2) = a - b(q_1 + q_2)\). Their costs are \(c\) per unit sold.



\textbf{(a)} Argue that successive elimination of strictly dominated strategies yields a unique prediction in this game.



\textbf{(b)} Would this be true if there were three firms instead of two?

\end{quote}

